% title: 단진자
% description: 길이 L인 줄에 매달려 연직선과 각 θ를 이룬 추에 작용하는 장력 T, 중력 mg와 그 성분 mg sin θ(복원력 방향), mg cos θ(줄 방향)
% packages:
% libraries: arrows.meta, calc, patterns
\begin{tikzpicture}[
    force/.style={very thick, -{Stealth[length=3mm]}},
    comp/.style={thin, dashed, -{Stealth[length=2mm]}}]
  \def\ang{30}    % 연직선과 줄 사이의 각 θ (도)
  \def\len{4}     % 줄의 길이 L
  \def\mg{2}      % 중력 화살표 길이
  \def\Tlen{1.6}  % 장력 화살표 길이
  \coordinate (P) at (0,0);                % 매단 점
  \coordinate (B) at (-90+\ang:\len);      % 추의 위치
  % 천장, 연직선(평형 위치 방향)
  \fill[pattern=north east lines] (-1,0) rectangle (1,0.25);
  \draw[thick] (-1,0) -- (1,0);
  \draw[thin, dashed] (P) -- (0,{-\len*cos(\ang)});
  % 줄과 각
  \draw[thick] (P) -- (B) node[pos=0.35, right] {$L$};
  \draw (0,-1) arc[start angle=-90, end angle=-90+\ang, radius=1];
  \node at (-90+\ang/2:1.3) {$\theta$};
  % 힘: 장력은 줄을 따라 매단 점 쪽, 중력은 아래
  \draw[force] (B) -- ++(90+\ang:\Tlen) node[left] {$\vec{T}$};
  \draw[force] (B) -- ++(0,-\mg) coordinate (W) node[below] {$m\vec{g}$};
  % 중력의 성분 (점선): 줄 방향 mg cos θ, 접선 방향 mg sin θ
  \draw[comp] (B) -- ++(-90+\ang:{\mg*cos(\ang)}) coordinate (Wr) node[right] {$mg\cos\theta$};
  \draw[comp] (B) -- ++(180+\ang:{\mg*sin(\ang)}) coordinate (Wt) node[left] {$mg\sin\theta$};
  \draw[very thin, dashed] (Wr) -- (W) -- (Wt);
  \draw ($(B)+(0,-0.7)$) arc[start angle=-90, end angle=-90+\ang, radius=0.7];
  \node at ($(B)+(-90+\ang/2:0.95)$) {$\theta$};
  \fill[blue!60!black] (B) circle[radius=0.18];
\end{tikzpicture}
